\documentclass[paper=a4, fontsize=11pt]{jlreq} % 日本語用の標準的なクラス

% --- パッケージ群 ---
\usepackage[utf8]{inputenc}
\usepackage[T1]{fontenc}
\usepackage{amsmath, amssymb, amsthm} % 数学の基本
\usepackage{physics} % \ket{0}, \pdv{x} などが簡単に打てる
\usepackage{bm} % 太字ベクトル \bm{v}
\usepackage{graphicx}
\usepackage{hyperref} % リンク・目次用
\usepackage{cite} % 引用用

% --- 定理環境の設定 (OMCの問題風に整理できる) ---
\theoremstyle{definition}
\newtheorem{theorem}{定理}[section]
\newtheorem{definition}[theorem]{定義}
\newtheorem{proposition}[theorem]{命題}
\newtheorem{example}[theorem]{例}
\newtheorem{remark}[theorem]{注釈}
\newtheorem{problem}[theorem]{問題} % 自作問題用

% --- 独自のショートカット (CFT/パンルヴェ用) ---
\newcommand{\Virexp}[2]{[L_{#1}, L_{#2}]} % ヴィラソロの交換関係
\newcommand{\taufunc}{\tau(z)} % タウ関数

\title{パンルヴェ方程式に関する勉強ノート}
\author{松浦光佑}
\date{\today}

\begin{document}

\maketitle
\tableofcontents % 自動で目次を作成

\section{はじめに}
本ノートは、岡本『パンルヴェ方程式』に関する学習記録である。
日々の学習を「AC（Accepted）」として記録し、GitHubでの進捗可視化を目的とする。

\section{パンルヴェ方程式の基礎}

    $x=\xi$をフックス型微分方程式
    \begin{equation}\label{eq:Fuchs}
    \frac{d^{2}y}{dx^{2}} + p_{1}(x)\frac{dy}{dx} + p_{2}(x)y = 0   
    \end{equation}
    の特異点とする。

\begin{proposition}[]
    % --- ここに数式 ---
    適当な2次の行列関数$\Phi_{\xi}(x)$を選んで
    \[
    \phi(x) = \vec{z}(x)\Phi_{\xi}(x)
    \]
    で定まる$\vec{z}(x)$が$x=\xi$のまわりで1価になるようにできる。もし、$\Gamma$が対角化可能ならば、$\Phi_{\xi}(x)$は対角行列になるように選べる。\\
    また、1価性により$x=\xi$のまわりで
    \begin{equation}\label{eq:power_series_z}
    \vec{z}(x) = \sum_{j=-\infty}^{\infty}\vec{a_{j}}u^{j}
    \end{equation}
    と表せる。
    \textbf{証明のスケッチ:}
    
    \begin{itemize}
        \item[\textbf{Insight:}] 通常、特性指数の差が整数の場合、解の片方に $\log$（対数項）が現れる。しかし、係数が特定の条件を満たして「$\log$ の項が完全に消滅する」特殊なケースがあり、それが非対数的特異点である。「みかけの特異点」の周りで解が分岐しないためにはこの性質が必須であり、この「対数項が消える条件」を計算することが、後のハミルトニアン導出の鍵を握る。
        \item[\textbf{Link:}] 岡本2.2節の命題2.5/p.59を参照。
    \end{itemize}
\end{proposition}

\begin{definition}[確定特異点]
    % --- ここに数式 ---
    \eqref{eq:power_series_z}において、ある整数$j_{0}$があって、$j\leq j_{0}$となるすべてのjに対し、
    \[
    \vec{a_{j}} = \vec{0}
    \]
    となる時$\xi$は\textbf{確定特異点}という。確定特異点でない時\textbf{不確定特異点}という。
    \begin{itemize}
        \item[\textbf{Insight:}] 特異点周りの級数展開において、「負のべき乗の項が無限に続かず、どこかで打ち止めになる（有限個しかない）」状態を表している。複素関数論における「極（pole）」に相当する、解析的に扱いやすい「おとなしい特異点」の定義である。パンルヴェ方程式の根幹である「動く特異点が極のみである（パンルヴェ性）」という性質を語る上で、この「確定（極）か、不確定（真性特異点など）か」の厳密な線引きがすべての土台となる。
        \item[\textbf{Link:}] 岡本2.2節の定義2.5/p.59を参照。
    \end{itemize}
\end{definition}

\begin{definition}[非対数的特異点]
    % --- ここに数式 ---
    $x=\xi$の特性指数$s_{1},s_{2}$が整数差のとき$x=\xi$は\textbf{非対数的特異点}と呼ばれる。
    \begin{itemize}
        \item[\textbf{Insight:}] 通常、特性指数の差が整数の場合、解の片方に $\log$（対数項）が現れる。しかし、係数が特定の条件を満たして「$\log$ の項が完全に消滅する」特殊なケースがあり、それが非対数的特異点である。「みかけの特異点」の周りで解が分岐しないためにはこの性質が必須であり、この「対数項が消える条件」を計算することが、後のハミルトニアン導出の鍵を握る。
        \item[\textbf{Link:}] 岡本2.2節の(2.20)/p.64を参照。
    \end{itemize}
\end{definition}

\begin{definition}[リーマン図式]
    % --- ここに数式 ---
    \[
    \left\{
    \begin{matrix}
        x = \xi_1 & x = \xi_2 & \cdots & x = \xi_m & x = \infty \\
        s_1^{(1)} & s_1^{(2)} & \cdots & s_1^{(m)} & s_1^{(\infty)} \\
        s_2^{(1)} & s_2^{(2)} & \cdots & s_2^{(m)} & s_2^{(\infty)} \\
        \vdots    & \vdots    &        & \vdots    & \vdots \\
        s_n^{(1)} & s_n^{(2)} & \cdots & s_n^{(m)} & s_n^{(\infty)}
    \end{matrix}
    \right\}
    \]
    \begin{itemize}
        \item[\textbf{Type:}] 確定特異点$\xi_{i}$と各確定特異点の特性指数$s_j^{(i)}$を書き並べたもの。
        \item[\textbf{Insight:}] この表の数値（局所的な性質）を\textbf{一切変えずに}、特異点の位置 $\xi_i$ だけを時間発展させた時に、裏で満たされるべき非線形の方程式が「ガルニエ系」。つまり、モノドロミー保存変形において\textbf{「絶対に定数として保存したいターゲット」}がこれ。
        \item[\textbf{Link:}] 岡本2.4節の(2.48)/p.81を参照。
    \end{itemize}
\end{definition}

\begin{definition}[みかけの特異点]
  % --- ここに数式 ---
    $\Xi$の元を確定特異点としてもつ微分方程式を考える。
    このときモノドロミー$\hat{\rho}$が等しく、さらに$\Xi\cup\Lambda$の元を確定特異点としてもつフックス型微分方程式が存在したとする。
    このとき、以下の2条件を満たす$\Lambda$の各元を\textbf{みかけの特異点}という。
    \begin{itemize}
        \item[(1)] $x=\lambda$における特性指数は全て整数。
        \item[(2)] $x=\lambda$は非対数的特異点
    \end{itemize}
    \begin{itemize}
        \item[\textbf{Insight:}] 方程式の係数には特異性が現れるが、解自体はそこで分岐を持たず正則（局所モノドロミーが単位行列）になる「偽物の特異点」。パンルヴェ方程式やガルニエ系の理論において最大のポイントは、**「このみかけの特異点の位置 $\lambda$ が、ハミルトン系における正準変数（未知関数 $q$）の役割を果たす」**という点にある。
        \item[\textbf{Link:}] 岡本2.5節/p.86を参照。
    \end{itemize}
\end{definition}



\begin{definition}[モノドロミー保存変形]
  % --- ここに数式 ---
    \begin{equation}\label{eq:before}
    \frac{d^{n}y}{dx^{n}}+\sum_{j=1}^{n} p_{j}(x)\frac{d^{n-1}y}{dx^{n-1}} = 0   
    \end{equation}
    という以下の(P1)-(P3)を満たすフックス型微分方程式を考える。
    \begin{itemize}
        \item[(P1)] $x=\xi_{1},...,\xi_{n},\xi_{n+1}=\infty$を確定特異点、$x=\lambda_{1},...,\lambda_{g}$をみかけの特異点としてもつ。
        \item[(P2)] 各$x=\xi_{i}$における特性指数のどの2つも、差が整数ではない。
        \item[(P3)] モノドロミー$\hat{\rho}$は既約である。
    \end{itemize}
    $p_{j}(x)$が$t$に解析的に依存しているとする。
    \begin{equation}\label{eq:after}
    \frac{d^{n}y}{dx^{n}}+\sum_{j=1}^{n} p_{j}(x,t)\frac{d^{n-1}y}{dx^{n-1}} = 0   
    \end{equation}
    が以下の条件(H)を満たす時、\eqref{eq:after}を\eqref{eq:before}の\textbf{モノドロミー保存変形}または\textbf{ホロノミック変形}という。
    \begin{itemize}
        \item[\textbf{Type:}] \eqref{eq:before}は$\mathbb{P}^{1}(\mathbb{C})$上定義されている。$U\subset\mathbb{C}^{L}$は領域。$t=(t_{1},...,t_{L})\in U$
        \item[\textbf{Insight:}] 方程式のパラメータ $t$ を動かしても、基本群の表現（モノドロミー）が変化しないような変形。Schlesinger系やPainlevé方程式の導出における出発点となる。
        \item[\textbf{Link:}] 岡本2.8節の定義2.14/p.105を参照。
    \end{itemize}
\end{definition}





% --- 今後の課題（TODOリスト） ---
\section{Next Step / To-Do}
\begin{itemize}
    \item [-] 確定特異点とは？
    \item [-] 特性指数とは？
    \item [-] モノドロミーが既約とはどういうこと？
    \item [-] モノドロミーとは?
\end{itemize}

\end{document}